% ===== 原文 pp.4–7 =====
% tr-06 片段 a：§2 组织纲要与阅读指南（原文 pp.4–7；p.8 起为 tr-01 的 §3）
% 起点：p.4 首段 "For the intersection number generating function …"（§1 末段，归本块写）；
% 终点：p.7 末 "… Expected values of geometric invariants [Mir10]."。
% 本块无编号陈述、无编号公式、无图。p.5/p.6 的 "Lemma."/"Theorem." 是原文引用
% 他人文献（[Mir07a]、[TWZ06]）的不编号陈述段，按正文本体排版，不建定理环境。

% [原文 p.4]
对相交数生成函数
\[
\mathbf{F}(\lambda,t_0,t_1,\ldots)
=\sum_{g=0}^{\infty}\lambda^{2g-2}\sum_{\{d_j\}}
\Big\langle\prod_{j=1}^{\infty}\tau_{d_j}\Big\rangle_g
\prod_{r\ge 0}t_r^{n_r}/n_r!,
\]
其中 $n_r=\#\{j\mid d_j=r\}$，且
\[
\langle\tau_{d_1}\cdots\tau_{d_n}\rangle_g
=\int_{\Mgnb}\prod_{j=1}^{n}\psi_j^{d_j},
\]
当 $\sum_{j=1}^n d_j=3g-3+n$ 时成立、而在其他情形下乘积 $\langle\tau_*\rangle$ 取零，
则指数 $\ee^{\mathbf{F}}$ 满足 Virasoro 代数约束，[Mir07b, 定理 6.1]。

% [Tutor Note: p.4 的生成函数把全部相交数打包成一条形式幂级数。$\sum_{\{d_j\}}$ 实际是对
%  有限支撑的多重指标求和；同一 $r$ 的 $n_r$ 个 $\tau_r$ 不可区分，故分母为 $n_r!$，这是标准
%  的指数型生成函数写法。可验证锚点：取 $g=0,n=3$、$d_1=d_2=d_3=0$，则 $n_0=3$，该项为
%  $\lambda^{-2}t_0^3/3!\cdot\langle\tau_0^3\rangle_0=\lambda^{-2}t_0^3/6$，其中
%  $\langle\tau_0^3\rangle_0=\int_{\Mbar_{0,3}}1=1$（$\Mbar_{0,3}$ 为一点）；这也说明为何要
%  约定 $\langle\tau_*\rangle$ 在 $\sum_j d_j\ne 3g-3+n$ 时取零。]

\section{组织纲要与阅读指南}

本纲要结合了 [Mir07a, Mir07b] 与 [Wlp10, 第 9 章] 的阐述。各讲的内容在下一节给出。

\begin{itemize}[leftmargin=1.6em,itemsep=2pt,topsep=2pt]
\item Teichmüller 空间、模空间、映射类群与辛几何。
  \begin{itemize}[leftmargin=1.6em,itemsep=2pt,topsep=2pt]
  \item Teichmüller 空间 $\Tgn$ 与模空间 $\Mgn$；给定长度的测地边界的带边双曲曲面的
    Teichmüller 空间 $\Tg(L_1,\ldots,L_n)$ 与模空间 $\Mg(L_1,\ldots,L_n)$；增广
    Teichmüller 空间与 Deligne--Mumford 型紧化；Dehn 扭转与映射类群 $\MCG$ 的作用。
  \item WP 辛几何，[Wlp10]。
    \begin{itemize}[leftmargin=1.6em,itemsep=2pt,topsep=2pt]
    \item 辛形式 $\omega=2\omega_{\WP}$（即 Kähler 形式）及其归一化。
    \item 测地长度函数 $\ell_\alpha$、Fenchel--Nielsen 无穷小扭转形变 $t_\alpha$，
      以及对偶公式 $\omega(\ ,t_\alpha)=\dd\ell_\alpha$。
    \item Teichmüller 空间的 Fenchel--Nielsen（FN）扭转--长度坐标 $(\tau_j,\ell_j)$，
      以及公式 $\omega=\sum_j\dd\ell_j\wedge\dd\tau_j$。
    \end{itemize}
  \item 中间模空间（intermediate moduli space）$\Mgn^\gamma$，其元素为对 $(R,\gamma)$
    ——一个曲面与一条加权多重曲线（weighted multicurve）$\gamma=\sum_j c_j\gamma_j$。
    \begin{itemize}[leftmargin=1.6em,itemsep=2pt,topsep=2pt]
    \item 覆盖塔（covering tower）
      \[
      \Tgn\longrightarrow\Mgn^\gamma\longrightarrow\Mgn.
      \]
    \item 加权多重曲线的稳定化子子群 $\Stab(\gamma)\subset\MCG$。覆盖塔的 $\MCG$
      覆叠陪集（deck cosets）。

% [Tutor Note: 读"覆盖塔"与"MCG 覆叠陪集"要抓住三层空间的关系：$\Tgn$ 是带标记的
%  Teichmüller 空间，$\Mgn$ 是模空间（即 $\Tgn/\MCG$），中间的 $\Mgn^\gamma$ 记录
%  "曲面 + 加权多重曲线 $\gamma$"，可理解为 $\Tgn$ 模去稳定化子 $\Stab(\gamma)$；因此
%  第二层到第三层之间覆叠的变换群正是 $\MCG/\Stab(\gamma)$，其元素即原文所说的 MCG deck
%  cosets。结构锚点：p.20--22 的余体积公式证明正是沿这条塔把模空间积分逐步展开为
%  $\Tgn/\Stab(\gamma)$ 上的积分（见 [Mir07a, 定理 7.1]）。]

% [原文 p.5]
    \item $\Mgn^\gamma$ 与 $\Mg(L)$ 上的辛结构。
    \end{itemize}
  \item 总长度 $\ell=\sum_j c_j\ell_{\gamma_j}$ 在 $\Tgn$ 与 $\Mgn^\gamma$ 中的水平集。

    \textbf{引理.}\ [Mir07a, 引理 7.2]。为体积递推与辛约化作准备。$\Mgn^\gamma$ 的一个
    有限覆盖是辛平面与低维模空间的纤维积。
  \end{itemize}

\item McShane--Mirzakhani 长度恒等式。
  \begin{itemize}[leftmargin=1.6em,itemsep=2pt,topsep=2pt]
  \item 集合 $\cB$ 由端点在边界上、相对于边界取同伦类的简单弧组成。与边界正交的测地弧的
    分类：在每个端点处与边界正交的简单测地线 $\Longleftrightarrow$ 不相交的边界区间对
    $\Longleftrightarrow$ 裤子的线框（wire frames）。
  \item Birman--Series：简单测地线测度为零。
  \item 有理指数函数（rational exponential function）$H$；双曲三角函数
    $\mathcal{D}$ 与 $\mathcal{R}$；其间的关系。
  \item 长度恒等式。

    \textbf{定理.}\ [Mir07a, 定理 1.3 与 4.2] 及 [TWZ06, 定理 1.8]。设 $R$ 是边界为
    $\beta_j$、长度为 $L_j$ 的双曲曲面，
    \[
    L_1=\sum_{\alpha_1,\alpha_2}\mathcal{D}\big(L_1,\ell_{\alpha_1}(R),\ell_{\alpha_2}(R)\big)
    +\sum_{j=2}^{n}\sum_{\alpha}\mathcal{R}\big(L_1,L_j,\ell_\alpha(R)\big),
    \]
    其中第一个和式取遍所有使 $\beta_1,\alpha_1,\alpha_2$ 围成一个嵌入的裤子的简单闭
    测地线的无序对，而双重和式取遍所有使 $\beta_1,\beta_j,\alpha$ 围成一个嵌入的裤子的
    简单闭测地线。
  \item [Mir07a, 第 8 节]——认识并理解该恒等式是 MCG 基本域的光滑类比。约化到较小的
    MCG 群的作用。
  \end{itemize}

% [Tutor Note: 核对 p.5 长度恒等式的一个有效办法是取 $n=1$：此时第二项消失，恒等式退化为
%  $L_1=\sum_{\alpha_1,\alpha_2}\mathcal{D}(L_1,\ell_{\alpha_1},\ell_{\alpha_2})$，即
%  McShane 的原始恒等式 [McS98]（单边界情形）；一般 $n$ 情形的双求和项则给出与其余边界
%  $\beta_j$ 的耦合。另注意第一个和式按简单闭测地线的"无序对"求和：$(\alpha_1,\alpha_2)$
%  与 $(\alpha_2,\alpha_1)$ 只计一次，否则每条边界裤子会被重复计数，这是照抄或改写此式时
%  最容易出错的地方。]

\item Mirzakhani 体积递推。
  \begin{itemize}[leftmargin=1.6em,itemsep=2pt,topsep=2pt]
  \item 余体积公式——把模空间积分写成按长度水平集的积分——中间模空间 $\Mgn^\gamma$ 的
    作用——[Mir07a, 定理 7.1] 与 [Wlp10, 定理 9.5]。
  \item 对 McShane--Mirzakhani 恒等式的应用。
    \begin{itemize}[leftmargin=1.6em,itemsep=2pt,topsep=2pt]
    \item 低维模空间体积的积分 $\mathcal{A}^{connected}_{g,n}$、
      $\mathcal{A}^{disconnected}_{g,n}$ 与 $\mathcal{B}_{g,n}$。
    \item 相应的连通与不连通边界裤子构型（boundary pants configurations）。
% [原文 p.6]
    \item 结合长度恒等式与余体积公式以得到主要结果；见 [Wlp10, 第 91 页底部、第 92 页]。

      \textbf{定理.}\ [Mir07a, 第 5 节与定理 8.1]。对 $(g,n)\neq(1,1),(0,3)$，体积
      $V_{g,n}(L)$ 满足
      \[
      \frac{\partial}{\partial L_1}L_1V_{g,n}(L)
      =\mathcal{A}^{connected}_{g,n}(L)+\mathcal{A}^{disconnected}_{g,n}(L)
      +\mathcal{B}_{g,n}(L).
      \]
    \item 积分 $\mathcal{A}^{*}_{*}(L)$ 与 $\mathcal{B}_{*}(L)$ 是边界长度的多项式，
      其系数为 Riemann zeta 特殊值之和；这些系数是 $\pi$ 的幂的正有理倍数。
    \end{itemize}
  \end{itemize}

\item 辛约化与 Duistermaat--Heckman 定理。
  \begin{itemize}[leftmargin=1.6em,itemsep=2pt,topsep=2pt]
  \item Deligne--Mumford 紧化上的 WP kappa 方程 $\omega=2\pi^2\kappa_1$ [Wlp90]。
  \item 边界为带点的测地边界的双曲曲面的 Teichmüller 空间
    $\what{\cT}_g(L)=\what{\cT}_g(L_1,\ldots,L_n)$ 的几何与拓扑。
  \item 通过在近乎紧的裤子（almost tight pants）上求和，为 $\what{\cT}_g(L)$ 赋予辛结构。
    边界长度矩映射 $R\in\what{\cT}_g(L)\mapsto L^2/2\in\R^n_+$。把扭转边界点看作
    Hamilton 流。
  \item 辛约化，[FO06]。$\Tg(L)$ 作为约化空间 $\what{\cT}_g(L)/(S^1)^n$。
    $\what{\cT}_g(L)$ 作为 $\Tgn(L)$ 上的主 $(S^1)^n$ 丛，以及小 $L$ 等价
    $\Tg(L)\approx\Tgn(0)$，
    \[
    \begin{array}{ccc}
    (S^1)^n & \longrightarrow & \what{\cT}_g(L)\\[-4pt]
    & & \big\downarrow\\[-4pt]
    & & \Tg(L)\approx\Tgn(0)
    \end{array}
    \]
  \item 圆上的点、穿刺点处的余切 $\C$-线、圆丛与同伦结构群。tautological 余切线类
    $\psi$。
  \item $\MCG$ 等变映射与商。Deligne--Mumford 型紧化与有限覆盖。椭圆栈 2。
  \item Duistermaat--Heckman 正规形式，[Mir07b, 定理 3.2]、[MS08, 第 2.5 节] 与
    [Wlp10, 第 95 页]，
    \[
    2\omega_{\what{\mathcal{M}}_g(L)/(S^1)^n}\equiv 2\omega_{\Mgnb(0)}+\sum_j\frac{L_j^2}{2}\psi_j.
    \]
  \item 辛约化的推论，[Mir07b, 定理 4.4]。

    \textbf{定理.}\ 体积多项式 $V_{g,n}(L)$ 的系数是 $\Mgnb$ 的相交数，由下式给出
    \[
    C_\alpha=\frac{2^{\delta_{1g}\delta_{1n}}}{2^{|\alpha|}\alpha!\,(3g-3+n-|\alpha|)!}
    \int_{\Mgnb}\psi_1^{\alpha_1}\cdots\psi_n^{\alpha_n}\omega^{3g-3+n-|\alpha|},
    \]
% [原文 p.7]
    其中 $\psi_j$ 是沿第 $j$ 个穿刺点的余切线的 Chern 类，$\omega$ 是辛形式，而
    $\delta_{**}$ 是 Kronecker 指示 delta。
  \end{itemize}

\item 相交数的模式。
  \begin{itemize}[leftmargin=1.6em,itemsep=2pt,topsep=2pt]
  \item 一般相交数记号
    \[
    \langle\kappa_1^{d_0}\tau_{d_1}\cdots\tau_{d_n}\rangle_g
    =\int_{\Mgnb}\prod_{j=1}^{n}\psi_j^{d_j}\kappa_1^{d_0}
    \]
    以及体积多项式展开。
  \item 例：$V_{1,1}(L)=\frac{\pi^2}{6}+\frac{L^2}{24}$，
    $V_{0,4}(L_1,L_2,L_3,L_4)=(4\pi^2+L_1^2+L_2^2+L_3^2+L_4^2)/2$，以及
    $V_{2,1}(L)=\frac{1}{2211840}(L^2+4\pi^2)(L^2+12\pi^2)
    (5L^4+384\pi^2L^2+6960\pi^4)$。
  \item 配分函数 $\mathbf{F}=\sum_g\big\langle \ee^{\sum_j t_j\tau_j}\big\rangle_g$ 与
    Virasoro 约束微分算子 $\mathbf{L}_k$ [Mir07b]。配分函数
    $\mathbf{G}=\sum_g\big\langle \ee^{s\kappa_1+\sum_j t_j\tau_j}\big\rangle_g$ 与
    Virasoro 约束微分算子 $\mathbf{V}_k$ [MS08, 定理 1.1]。
  \item Do 的\emph{移除边界}（remove a boundary）关系
    \[
    \frac{\partial V_{g,n+1}}{\partial L_{n+1}}(L,2\pi\ii)=2\pi\ii(2g-2+n)V_{g,n}(L),
    \qquad [Do08, DN09].
    \]
  \item Manin--Zograf 体积生成函数 [MZ00]。穿刺点渐近——对正常数 $c$、$a_g$，固定亏格、
    $n$ 很大时，有
    \[
    V_{g,n}=c^n n!\,n^{(5g-7)/2}(a_g+O(1/n)),\qquad [MZ00].
    \]
  \item Schumacher--Trapani 亏格渐近——固定 $n$ 时存在正常数，使得
    \[
    c_1^g(2g)!<V_{g,n}<c_2^g(2g)!,\qquad [Gru01, ST01].
    \]
  \item Zograf 猜想
    \[
    V_{g,n}=(4\pi^2)^{2g+n-3}(2g+n-3)!\frac{1}{\sqrt{g\pi}}
    \Big(1+\frac{c_n}{g}+O(1/g^2)\Big),
    \]
    对固定 $n$、$g$ 趋于无穷成立，[Zog08]。几何不变量的期望值 [Mir10]。
  \end{itemize}
\end{itemize}

% [Tutor Note: p.6 的 "The elliptic stack 2." 已用 1000 dpi 裁图复核：字形确实就是数字
%  "2"，既非上标也非脚注标记，该页也没有任何脚注。按上下文（在 Deligne--Mumford 型紧化
%  与有限覆盖之后），这里很可能是原稿丢失了一个数学记号（例如椭圆栈 $\Mbar_{1,1}$ 的
%  写法），但底本此处只能照录为"椭圆栈 2"，读者留意这处存疑。]
%
% [Tutor Note: Zograf 渐近式的分母是 $\sqrt{g\pi}$（根号横线罩住 $g\pi$ 两个字符），
%  不是 $\sqrt{g}\,\pi$；已用 600 dpi 裁图确认。量级自检：取 $g=2,n=0$，该式给出
%  $(4\pi^2)^{1}\cdot 1!/\sqrt{2\pi}\approx 15.75$，而 p.35 的精确值
%  $V_{2,0}=43\pi^6/2160\approx 19.14$，比值 $\approx1.215$ 对应 $1+c_0/2$，即
%  $c_0\approx0.43$，量级合理；若误读成分母 $\sqrt{g}\,\pi$，同一估值只有 $\approx 8.89$，
%  与精确值相差两倍以上，错误立刻暴露。]
